\documentclass[12pt]{article}
\usepackage[T1]{fontenc}
\usepackage[utf8]{inputenc}
\usepackage{lmodern}
\usepackage{textcomp}
\usepackage{enumitem}
\usepackage{geometry}
\geometry{margin=1in}
\begin{document}
\begin{center}
Answer Key - Science - K-12 Level Assessment 22 \
\small (Answers are ordered exactly as in the question paper.)
\end{center}

\section*{Multiple Choice}
\begin{enumerate}[label=\arabic*.]
\item \textbf{Answer:} D
\\ \textit{Options:} A) heat; B) chemical; C) mechanical; D) electromagnetic
\\ \textit{Note:} The correct answer is electromagnetic because when electrical energy flows through the filament in a light bulb, it excites the atoms, causing them to emit energy in the form of electromagnetic radiation, which includes visible light.
\item \textbf{Answer:} D
\\ \textit{Options:} A) tensile; B) shear; C) torsion; D) compression
\\ \textit{Note:} The primary type of stress that a fixed vise exerts on a piece of wood is compression because the vise applies inward pressure on the wood, forcing its fibers to be pushed together and creating a compressive force.
\item \textbf{Answer:} D
\\ \textit{Options:} A) because it forms where fresh and salt water meet; B) because it receives water directly from precipitation; C) because it rises to the surface near the ocean; D) because pollutants are filtered by rock and soil deep within Earth
\\ \textit{Note:} Water from an aquifer is more likely to be cleaner than water from other sources because as it passes through layers of rock and soil, these natural materials filter out pollutants, resulting in purer water.
\item \textbf{Answer:} B
\\ \textit{Options:} A) thermal; B) solar; C) mechanical; D) chemical
\\ \textit{Note:} Solar energy, in the form of electromagnetic radiation, can travel through a vacuum because it does not require a medium to propagate, unlike sound or mechanical waves.
\item \textbf{Answer:} A
\\ \textit{Options:} A) Rock cools quickly from melted rock.; B) Rock is recrystallized by extreme pressure.; C) Rock solidifies slowly deep underground.; D) Rock is formed from deposited sediment.
\\ \textit{Note:} The correct answer is based on the fact that when molten rock, or magma, erupts from a volcano and comes into contact with the cooler surface environment, it rapidly loses heat and solidifies, leading to the formation of new surface rock.
\end{enumerate}

\section*{True / False}
\begin{enumerate}[label=\arabic*.]
\setcounter{enumi}{5}
\item \textbf{Answer:} False
\\ \textit{Options:} A) True; B) False
\\ \textit{Note:} A toy truck will roll slower on a sandy surface than on a smooth surface because the sand creates more friction and resistance against the wheels, hindering its movement.
\item \textbf{Answer:} False
\\ \textit{Options:} A) True; B) False
\\ \textit{Note:} The primary energy source for deep ocean currents is the wind and the Earth's rotation rather than the tidal effects of the Moon.
\item \textbf{Answer:} False
\\ \textit{Options:} A) True; B) False
\\ \textit{Note:} A thermal insulator is a material that resists the transfer of heat, while specific heat capacity refers to the amount of heat energy required to raise the temperature of a material, so a high specific heat capacity does not necessarily indicate insulation properties.
\item \textbf{Answer:} True
\\ \textit{Options:} A) True; B) False
\\ \textit{Note:} Matter is defined as anything that has mass and occupies space, thereby possessing both mass and volume as fundamental properties of all substances.
\item \textbf{Answer:} True
\\ \textit{Options:} A) True; B) False
\\ \textit{Note:} The statement is true because an object in motion will remain in motion at a constant velocity unless acted upon by an unbalanced force, so if a book is moving horizontally, it indicates that a force is acting on it.
\end{enumerate}

\section*{Long Form Response}
\begin{enumerate}[label=\arabic*.]
\setcounter{enumi}{10}
\item \textbf{Answer:} Chromosomes play a crucial role in human reproduction by carrying the genetic information that determines the traits of an offspring. When male and female gametes, which are sperm and egg cells respectively, combine during fertilization, they each contribute half of the necessary chromosomes. In humans, this means that the sperm carries 23 chromosomes, and the egg also carries 23 chromosomes, resulting in a total of 46 chromosomes in the fertilized egg. This combination of genetic material from both parents not only forms a new individual but also ensures genetic diversity, as the offspring inherits a unique set of genes from each parent. Thus, the union of male and female gametes is essential for the development of a new organism.
\\ \textit{Note:} correct because it accurately describes how chromosomes, contributed equally by male and female gametes during fertilization, combine to form a complete set of 46 chromosomes in the offspring, which carries the genetic information essential for determining traits.
\item \textbf{Answer:} Wind drag plays a significant role in shaping ocean surface currents by transferring energy from the wind to the water. As the wind blows across the ocean's surface, it creates friction, which causes the water to move and form currents. These currents not only transport warm and cold water across vast distances but also influence weather patterns and climate. Additionally, the movement of these currents affects marine ecosystems by distributing nutrients and supporting diverse marine life. Therefore, understanding wind drag is essential for grasping how ocean currents are formed, their movement, and their broader impact on the environment.
\\ \textit{Note:} correct because it effectively explains how wind drag initiates and sustains ocean surface currents, highlighting their role in energy transfer, climate influence, and nutrient distribution which are crucial for marine ecosystems.
\end{enumerate}

\vfill
\noindent\textit{End of Answer Key}
\end{document}